\documentclass[aspectratio=169]{beamer}
\usetheme{metropolis}
\usepackage{graphicx}
\usepackage{tikz}
\usepackage{amsmath}
\usepackage{listings}
\usepackage{xcolor}
\usepackage{booktabs}
\usepackage{emoji}
\usetikzlibrary{shapes,arrows,positioning,calc}

% Code listing setup
\lstset{
  basicstyle=\ttfamily\small,
  keywordstyle=\color{blue},
  commentstyle=\color{gray},
  stringstyle=\color{red},
  showstringspaces=false,
  breaklines=true,
  frame=single,
  language=Python
}

\title{Lecture 15: BERT Deep Dive}
\subtitle{Week 6, Lecture 1 - Bidirectional Encoder Representations}
\author{LLM Course}
\date{Winter 2026}

\begin{document}

% ============================================================
% TITLE SLIDE
% ============================================================
\begin{frame}
\titlepage
\end{frame}

% ============================================================
% OVERVIEW
% ============================================================
\begin{frame}{Today's Agenda 📋}
\large
\begin{enumerate}
    \item 🎯 \textbf{BERT Introduction}: What makes it special?
    \item 🎭 \textbf{Masked Language Modeling}: The key training objective
    \item 🏗️ \textbf{BERT Architecture}: Model sizes and specifications
    \item 📊 \textbf{Pre-training \& Fine-tuning}: The two-stage paradigm
    \item 🔬 \textbf{Contextual Embeddings}: Seeing polysemy in action
    \item 💻 \textbf{Using BERT}: Practical code examples
\end{enumerate}

\vspace{1em}
\textit{Goal: Deep understanding of BERT and how it revolutionized NLP}
\end{frame}

% ============================================================
% BERT INTRO
% ============================================================
\begin{frame}{BERT: Bidirectional Encoder Representations 🎯}
\begin{center}
\LARGE
\textbf{BERT = Encoder-only Transformer}
\end{center}

\vspace{1em}

\textbf{Key Innovation: Masked Language Modeling (MLM)}

\vspace{1em}

\begin{columns}[T]
\begin{column}{0.48\textwidth}
\textbf{Traditional Language Models:}
\begin{itemize}
    \item Left-to-right (GPT)
    \item Or right-to-left
    \item \alert{Unidirectional context}
    \item Can't see full picture
\end{itemize}

\vspace{1em}

\textbf{Example:}
\begin{itemize}
    \item "The cat sat on the \_\_\_"
    \item Only sees left context
\end{itemize}
\end{column}

\pause

\begin{column}{0.48\textwidth}
\textbf{BERT (MLM):}
\begin{itemize}
    \item Mask random tokens
    \item Predict from \alert{bidirectional context}
    \item Sees both left \& right
    \item Deeper understanding
\end{itemize}

\vspace{1em}

\textbf{Example:}
\begin{itemize}
    \item "The cat [MASK] on the mat"
    \item Sees: "The cat" AND "on the mat"
    \item Predicts: "sat"
\end{itemize}
\end{column}
\end{columns}

\vspace{1em}
\small
\textit{Reference: Devlin et al. (2019) - "BERT: Pre-training of Deep Bidirectional Transformers"}
\end{frame}

% ============================================================
% WHY BERT
% ============================================================
\begin{frame}{Why BERT Was Revolutionary 🚀}
\textbf{Before BERT (2018):}
\begin{itemize}
    \item Feature-based approaches (use Word2Vec/GloVe as features)
    \item Task-specific architectures
    \item Limited transfer learning
    \item Unidirectional or shallow bidirectional models
\end{itemize}

\vspace{1em}

\textbf{BERT's Contributions:}
\begin{enumerate}
    \item \textbf{Deep Bidirectionality}
    \begin{itemize}
        \item True bidirectional context at every layer
        \item Not just concatenating left-to-right and right-to-left
    \end{itemize}

    \item \textbf{Pre-train + Fine-tune Paradigm}
    \begin{itemize}
        \item Single pre-trained model for all tasks
        \item Fine-tune with minimal architecture changes
        \item Democratized NLP (no need to train from scratch!)
    \end{itemize}

    \item \textbf{State-of-the-Art Results}
    \begin{itemize}
        \item Beat previous best on 11 NLP tasks
        \item Large performance gains (sometimes 10+ points!)
        \item Showed power of pre-training
    \end{itemize}
\end{enumerate}
\end{frame}

% ============================================================
% MLM TRAINING
% ============================================================
\begin{frame}{Masked Language Modeling (MLM) 🎭}
\textbf{BERT's Pre-training Objective}

\vspace{1em}

\textbf{Training Procedure:}
\begin{enumerate}
    \item Take a sentence
    \item Randomly mask 15\% of tokens
    \item Of the masked tokens:
    \begin{itemize}
        \item 80\%: Replace with [MASK]
        \item 10\%: Replace with random word
        \item 10\%: Keep unchanged
    \end{itemize}
    \item Predict the original tokens
\end{enumerate}

\vspace{1em}

\begin{exampleblock}{Example}
\textbf{Original:} "My dog is hairy"

\textbf{Masked:} "My dog is [MASK]"

\textbf{Labels:} [-, -, -, hairy]

\textbf{Prediction:} Model predicts "hairy" using bidirectional context
\end{exampleblock}

\vspace{1em}

\textbf{Why the 80/10/10 split?}
\begin{itemize}
    \item Prevents overfitting to [MASK] token
    \item Forces model to maintain representations for all tokens
\end{itemize}
\end{frame}

% ============================================================
% MLM EXAMPLE
% ============================================================
\begin{frame}{MLM Example: Step by Step 📝}
\textbf{Sentence:} "The quick brown fox jumps over the lazy dog"

\vspace{1em}

\textbf{Step 1: Random Masking (15\% of tokens)}

\begin{center}
\begin{tikzpicture}[scale=0.8]
    \foreach \i/\word in {1/The, 2/quick, 3/brown, 4/fox, 5/jumps, 6/over, 7/the, 8/lazy, 9/dog} {
        \node at (\i*1.3, 0) {\word};
    }

    % Highlight masked tokens
    \node[rectangle,draw,fill=orange!30] at (2.6, 0) {quick};
    \node[rectangle,draw,fill=red!30] at (6.5, 0) {over};

    \node at (6, -1) {Select 2 tokens for masking (15\% of 9)};
\end{tikzpicture}
\end{center}

\vspace{1em}

\textbf{Step 2: Apply Masking Strategy}
\begin{itemize}
    \item "quick" → 80\% chance → [MASK]
    \item "over" → 10\% chance → "under" (random word)
\end{itemize}

\vspace{0.5em}

\textbf{Input to BERT:} "The [MASK] brown fox jumps under the lazy dog"

\textbf{Prediction targets:} quick, over

\vspace{1em}

\textbf{Model learns:}
\begin{itemize}
    \item "quick" from context: "The \_\_\_ brown" (adjective before noun)
    \item "over" from context: "jumps \_\_\_ the" (preposition in this context)
\end{itemize}
\end{frame}

% ============================================================
% NSP
% ============================================================
\begin{frame}{Next Sentence Prediction (NSP) 🔗}
\textbf{BERT's second pre-training objective (debated usefulness)}

\vspace{1em}

\textbf{Task:} Given two sentences A and B, predict if B follows A

\vspace{1em}

\begin{exampleblock}{Positive Example (IsNext)}
\textbf{Sentence A:} "The man went to the store."

\textbf{Sentence B:} "He bought a gallon of milk."

\textbf{Label:} IsNext ✓
\end{exampleblock}

\vspace{0.5em}

\begin{alertblock}{Negative Example (NotNext)}
\textbf{Sentence A:} "The man went to the store."

\textbf{Sentence B:} "Penguins are flightless birds."

\textbf{Label:} NotNext ✗
\end{alertblock}

\vspace{1em}

\textbf{Implementation:}
\begin{itemize}
    \item Use [CLS] token representation for classification
    \item 50\% real pairs, 50\% random pairs
    \item Binary classification task
\end{itemize}

\vspace{0.5em}
\textit{Note: Later work (RoBERTa) showed NSP might not be necessary!}
\end{frame}

% ============================================================
% BERT ARCHITECTURE
% ============================================================
\begin{frame}{BERT Architecture Variants 📊}
\begin{center}
\begin{tabular}{lcccc}
\toprule
\textbf{Model} & \textbf{Layers} & \textbf{Hidden Size} & \textbf{Heads} & \textbf{Params} \\
\midrule
BERT-Base & 12 & 768 & 12 & 110M \\
BERT-Large & 24 & 1024 & 16 & 340M \\
\bottomrule
\end{tabular}
\end{center}

\vspace{1em}

\textbf{Architecture Details (BERT-Base):}
\begin{itemize}
    \item 12 transformer encoder layers
    \item 768-dimensional hidden states
    \item 12 attention heads per layer (64 dims each)
    \item 3072-dimensional feed-forward intermediate size (4x expansion)
    \item Maximum sequence length: 512 tokens
    \item Vocabulary size: 30,000 WordPiece tokens
\end{itemize}

\vspace{1em}

\textbf{Special Tokens:}
\begin{itemize}
    \item \textbf{[CLS]}: Classification token (first token, used for sequence-level tasks)
    \item \textbf{[SEP]}: Separator token (between sentences)
    \item \textbf{[MASK]}: Mask token (for MLM)
    \item \textbf{[PAD]}: Padding token (for variable-length sequences)
\end{itemize}
\end{frame}

% ============================================================
% BERT INPUT
% ============================================================
\begin{frame}{BERT Input Representation 🔤}
\textbf{Three types of embeddings are summed:}

\vspace{1em}

\begin{center}
\begin{tikzpicture}[scale=0.8]
    % Input text
    \node at (0, 5) {\textbf{Input:}};
    \node at (2, 5) {[CLS]};
    \node at (3.5, 5) {my};
    \node at (4.8, 5) {dog};
    \node at (6.1, 5) {is};
    \node at (7.5, 5) {cute};
    \node at (9, 5) {[SEP]};
    \node at (10.5, 5) {he};
    \node at (12, 5) {plays};

    % Token embeddings
    \node at (0, 3.5) {\small\textbf{Token:}};
    \foreach \x in {2, 3.5, 4.8, 6.1, 7.5, 9, 10.5, 12} {
        \node[rectangle,draw,fill=blue!20,minimum width=1cm,minimum height=0.6cm] at (\x, 3.5) {};
    }

    % Segment embeddings
    \node at (0, 2) {\small\textbf{Segment:}};
    \foreach \x in {2, 3.5, 4.8, 6.1, 7.5} {
        \node[rectangle,draw,fill=green!20,minimum width=1cm,minimum height=0.6cm] at (\x, 2) {A};
    }
    \foreach \x in {9, 10.5, 12} {
        \node[rectangle,draw,fill=orange!20,minimum width=1cm,minimum height=0.6cm] at (\x, 2) {B};
    }

    % Position embeddings
    \node at (0, 0.5) {\small\textbf{Position:}};
    \foreach \x/\pos in {2/0, 3.5/1, 4.8/2, 6.1/3, 7.5/4, 9/5, 10.5/6, 12/7} {
        \node[rectangle,draw,fill=purple!20,minimum width=1cm,minimum height=0.6cm] at (\x, 0.5) {\pos};
    }

    % Sum
    \draw[->, thick] (7, -0.5) -- (7, -1);
    \node at (7, -1.5) {\textbf{Sum all three embeddings}};

    \draw[->, thick] (7, -2) -- (7, -2.5);
    \node[rectangle,draw,fill=red!20,minimum width=3cm,minimum height=0.8cm] at (7, -3) {BERT Encoder};
\end{tikzpicture}
\end{center}

\vspace{0.5em}

\textbf{Three embedding types:}
\begin{enumerate}
    \item \textbf{Token Embeddings}: WordPiece vocabulary
    \item \textbf{Segment Embeddings}: Which sentence (A or B)?
    \item \textbf{Position Embeddings}: Learned position (0 to 511)
\end{enumerate}
\end{frame}

% ============================================================
% PRE-TRAINING
% ============================================================
\begin{frame}{BERT Pre-training 🏋️}
\textbf{Massive scale pre-training on unlabeled text}

\vspace{1em}

\textbf{Pre-training Data:}
\begin{itemize}
    \item \textbf{BooksCorpus}: 800M words (novels, fiction)
    \item \textbf{English Wikipedia}: 2,500M words
    \item Total: 3.3 billion words
    \item Diverse, high-quality text
\end{itemize}

\vspace{1em}

\textbf{Training Details:}
\begin{itemize}
    \item Batch size: 256 sequences (128,000 tokens)
    \item Training steps: 1M steps
    \item Optimization: Adam (lr=1e-4, warmup=10k steps)
    \item Hardware: 4-16 Cloud TPUs
    \item Training time: 4 days (BERT-Base), 4+ days (BERT-Large)
\end{itemize}

\vspace{1em}

\begin{block}{Key Insight}
Pre-training learns general language understanding that transfers to many downstream tasks!
\end{block}
\end{frame}

% ============================================================
% FINE-TUNING
% ============================================================
\begin{frame}{Fine-tuning BERT 🎓}
\textbf{Two-stage process: Pre-train then Fine-tune}

\vspace{1em}

\begin{center}
\begin{tikzpicture}[scale=0.8]
    % Pre-training
    \node[rectangle,draw,fill=blue!20,minimum width=4cm,minimum height=1.5cm] (pretrain) at (0, 3) {
        \begin{tabular}{c}
        \textbf{Pre-training} \\
        Unlabeled text \\
        MLM + NSP \\
        Days on TPUs
        \end{tabular}
    };

    \node[above=0.1cm of pretrain] {\textbf{Stage 1}};

    % Arrow
    \draw[->, very thick] (0, 1.5) -- (0, 0.5);

    % Fine-tuning
    \node[rectangle,draw,fill=green!20,minimum width=4cm,minimum height=1.5cm] (finetune) at (0, -1) {
        \begin{tabular}{c}
        \textbf{Fine-tuning} \\
        Labeled task data \\
        Task-specific head \\
        Minutes to hours
        \end{tabular}
    };

    \node[below=0.1cm of finetune] {\textbf{Stage 2}};

    % Task examples
    \node[right=2cm of pretrain] {
        \begin{tabular}{l}
        \textbf{Pre-train once} \\
        Use for all tasks
        \end{tabular}
    };

    \node[right=2cm of finetune] {
        \begin{tabular}{l}
        \textbf{Fine-tune per task:} \\
        • Classification \\
        • NER \\
        • Question Answering \\
        • etc.
        \end{tabular}
    };
\end{tikzpicture}
\end{center}

\vspace{1em}

\textbf{Benefits:}
\begin{itemize}
    \item Pre-training is expensive but done once
    \item Fine-tuning is cheap and fast
    \item Same pre-trained model for all downstream tasks
\end{itemize}
\end{frame}

% ============================================================
% FINE-TUNING TASKS
% ============================================================
\begin{frame}{Fine-tuning for Different Tasks 🎯}
\textbf{Minimal architecture changes needed!}

\vspace{1em}

\begin{enumerate}
    \item \textbf{Single Sentence Classification}
    \begin{itemize}
        \item Input: [CLS] sentence [SEP]
        \item Output: [CLS] representation → classifier
        \item Example: Sentiment analysis
    \end{itemize}

    \item \textbf{Sentence Pair Classification}
    \begin{itemize}
        \item Input: [CLS] sentence A [SEP] sentence B [SEP]
        \item Output: [CLS] representation → classifier
        \item Example: Natural Language Inference
    \end{itemize}

    \item \textbf{Question Answering}
    \begin{itemize}
        \item Input: [CLS] question [SEP] passage [SEP]
        \item Output: Token-level predictions for start/end positions
        \item Example: SQuAD
    \end{itemize}

    \item \textbf{Token Classification}
    \begin{itemize}
        \item Input: [CLS] sentence [SEP]
        \item Output: Each token representation → classifier
        \item Example: Named Entity Recognition
    \end{itemize}
\end{enumerate}
\end{frame}

% ============================================================
% CODE EXAMPLE
% ============================================================
\begin{frame}[fragile]{Fine-tuning BERT: Code Example 💻}
\textbf{Using HuggingFace Transformers}

\begin{lstlisting}[language=Python]
from transformers import BertForSequenceClassification, Trainer, TrainingArguments

# Load pre-trained BERT with classification head
model = BertForSequenceClassification.from_pretrained(
    'bert-base-uncased',
    num_labels=2  # Binary classification
)

# Define training arguments
training_args = TrainingArguments(
    output_dir='./results',
    num_train_epochs=3,
    per_device_train_batch_size=16,
    learning_rate=2e-5,
    warmup_steps=500,
)

# Train
trainer = Trainer(
    model=model,
    args=training_args,
    train_dataset=train_dataset,
    eval_dataset=eval_dataset,
)

trainer.train()
\end{lstlisting}

\vspace{0.5em}
\small
\textit{Reference: HuggingFace Course - Chapters 1.5, 7.3}
\end{frame}

% ============================================================
% CONTEXTUAL EMBEDDINGS EXAMPLE
% ============================================================
\begin{frame}[fragile]{Contextual Embeddings in Action 🔬}
\textbf{Remember "bank"? Let's see BERT handle it!}

\begin{lstlisting}[language=Python]
from transformers import BertTokenizer, BertModel
import torch

tokenizer = BertTokenizer.from_pretrained('bert-base-uncased')
model = BertModel.from_pretrained('bert-base-uncased')

# Two different contexts for "bank"
sent1 = "I deposited money at the bank"
sent2 = "We sat by the river bank"

# Get embeddings
def get_embedding(sentence, target_word):
    inputs = tokenizer(sentence, return_tensors='pt')
    outputs = model(**inputs)
    # Find position of target word
    tokens = tokenizer.tokenize(sentence)
    idx = tokens.index(target_word) + 1  # +1 for [CLS]
    return outputs.last_hidden_state[0, idx, :]

emb1 = get_embedding(sent1, "bank")  # Financial bank
emb2 = get_embedding(sent2, "bank")  # River bank

# Compare similarity
similarity = torch.cosine_similarity(emb1, emb2, dim=0)
print(f"Similarity: {similarity:.3f}")  # Low! (~0.3-0.5)
# Different contexts → Different embeddings!
\end{lstlisting}
\end{frame}

% ============================================================
% BERT RESULTS
% ============================================================
\begin{frame}{BERT's Impressive Results 📈}
\textbf{State-of-the-art on 11 NLP tasks when released (2018)}

\vspace{1em}

\begin{center}
\small
\begin{tabular}{lccc}
\toprule
\textbf{Task} & \textbf{Metric} & \textbf{Pre-BERT} & \textbf{BERT} \\
\midrule
SQuAD 1.1 (QA) & F1 & 84.1 & \textbf{90.9} \\
SQuAD 2.0 (QA) & F1 & 66.3 & \textbf{83.1} \\
MNLI (NLI) & Accuracy & 80.6 & \textbf{86.7} \\
SST-2 (Sentiment) & Accuracy & 93.2 & \textbf{94.9} \\
CoNLL-2003 (NER) & F1 & 92.6 & \textbf{92.8} \\
\bottomrule
\end{tabular}
\end{center}

\vspace{1em}

\textbf{Key Observations:}
\begin{itemize}
    \item Largest gains on tasks requiring understanding (QA, NLI)
    \item Improvements even on well-studied benchmarks
    \item BERT-Large generally better than BERT-Base
    \item Fine-tuning is simple but very effective
\end{itemize}

\vspace{1em}

\begin{block}{Impact}
BERT made pre-trained transformers the standard approach in NLP. Almost all subsequent models build on BERT's ideas!
\end{block}
\end{frame}

% ============================================================
% WHAT BERT LEARNS
% ============================================================
\begin{frame}{What Does BERT Learn? 🧠}
\textbf{Probing BERT's internal representations}

\vspace{1em}

\textbf{Research has shown BERT captures:}

\begin{enumerate}
    \item \textbf{Syntactic Information}
    \begin{itemize}
        \item Part-of-speech tags
        \item Constituent structure
        \item Dependency relations
        \item Lower layers encode more syntax
    \end{itemize}

    \item \textbf{Semantic Information}
    \begin{itemize}
        \item Word sense disambiguation
        \item Semantic roles
        \item Entity types
        \item Middle layers encode more semantics
    \end{itemize}

    \item \textbf{Pragmatic Information}
    \begin{itemize}
        \item Coreference resolution
        \item Discourse relations
        \item Higher layers encode more pragmatics
    \end{itemize}

    \item \textbf{World Knowledge}
    \begin{itemize}
        \item Factual knowledge (to some extent)
        \item Common sense reasoning (limited)
    \end{itemize}
\end{enumerate}

\vspace{0.5em}
\small
\textit{Reference: Tenney et al. (2019) - "BERT Rediscovers the Classical NLP Pipeline"}
\end{frame}

% ============================================================
% LAYER ANALYSIS
% ============================================================
\begin{frame}{BERT Layer Analysis 📊}
\textbf{Different layers capture different linguistic properties}

\vspace{1em}

\begin{center}
\begin{tikzpicture}[scale=0.9]
    % Layers
    \foreach \y/\layer in {0/Layer 1, 1/Layer 4, 2/Layer 8, 3/Layer 12} {
        \node[rectangle,draw,fill=blue!20,minimum width=3cm,minimum height=0.8cm] at (0, \y*1.5) {\layer};
    }

    % Properties
    \node[right=1cm of {(0,0)}] {\textbf{Lexical/Surface}};
    \node[right=1cm of {(0,1.5)}] {\textbf{Syntactic}};
    \node[right=1cm of {(0,3)}] {\textbf{Semantic}};
    \node[right=1cm of {(0,4.5)}] {\textbf{Task-Specific}};

    % Gradient
    \fill[blue!10] (6, 0) rectangle (8, 0.8);
    \fill[blue!30] (6, 0.8) rectangle (8, 2.3);
    \fill[blue!50] (6, 2.3) rectangle (8, 3.8);
    \fill[blue!70] (6, 3.8) rectangle (8, 5.3);

    \node[right=0.2cm of {(8, 2.65)}] {Abstraction};
    \draw[->, thick] (8.5, 0.4) -- (8.5, 5);
\end{tikzpicture}
\end{center}

\vspace{1em}

\textbf{Observations:}
\begin{itemize}
    \item Lower layers: surface features (word forms, POS)
    \item Middle layers: syntax (phrase structure, dependencies)
    \item Higher layers: semantics and task-specific features
    \item Hierarchical representation learning
\end{itemize}

\vspace{0.5em}
\textit{Similar to how CNNs learn in computer vision: edges → shapes → objects}
\end{frame}

% ============================================================
% DISCUSSION
% ============================================================
\begin{frame}{Discussion Questions 💭}
\begin{enumerate}
    \item \textbf{MLM vs. Autoregressive:}
    \begin{itemize}
        \item Why is MLM better for understanding tasks?
        \item Can BERT generate text like GPT?
        \item What are the trade-offs?
    \end{itemize}

    \vspace{0.5em}

    \item \textbf{The 80/10/10 Masking Strategy:}
    \begin{itemize}
        \item Why not just use 100\% [MASK]?
        \item What problem does the random replacement solve?
        \item Could we improve this strategy?
    \end{itemize}

    \vspace{0.5em}

    \item \textbf{Pre-training Data:}
    \begin{itemize}
        \item Why use books and Wikipedia?
        \item Would social media text work as well?
        \item How does data quality affect pre-training?
    \end{itemize}

    \vspace{0.5em}

    \item \textbf{Fine-tuning:}
    \begin{itemize}
        \item Why does fine-tuning work so well?
        \item When might fine-tuning fail?
        \item How much labeled data do we need?
    \end{itemize}
\end{enumerate}
\end{frame}

% ============================================================
% LOOKING AHEAD
% ============================================================
\begin{frame}{Looking Ahead 🔮}
\textbf{Today we learned:}
\begin{itemize}
    \item BERT architecture and innovations
    \item Masked Language Modeling
    \item Pre-training and fine-tuning
    \item Contextual embeddings
    \item What BERT learns
\end{itemize}

\vspace{1em}

\textbf{Next lecture (Lecture 16 - BERT Variants):}
\begin{itemize}
    \item \alert{RoBERTa}: Optimized BERT training
    \item \alert{ALBERT}: Parameter-efficient BERT
    \item \alert{DistilBERT}: Smaller, faster BERT
    \item \alert{Other variants}: ELECTRA, DeBERTa, and more
    \item Comparative analysis and when to use which
\end{itemize}

\vspace{1em}

\begin{center}
\Large
\textbf{BERT started a revolution in NLP! 🚀}
\end{center}
\end{frame}

% ============================================================
% SUMMARY
% ============================================================
\begin{frame}{Summary 🎯}
\textbf{Key Takeaways:}

\vspace{1em}

\begin{enumerate}
    \item \textbf{BERT = Encoder-only Transformer}
    \begin{itemize}
        \item Bidirectional self-attention
        \item Trained with Masked Language Modeling
    \end{itemize}

    \item \textbf{Pre-train + Fine-tune Paradigm}
    \begin{itemize}
        \item Expensive pre-training on unlabeled data (once)
        \item Cheap fine-tuning on task-specific data (per task)
    \end{itemize}

    \item \textbf{Contextual Embeddings}
    \begin{itemize}
        \item Different representations based on context
        \item Solves polysemy problem
    \end{itemize}

    \item \textbf{Hierarchical Learning}
    \begin{itemize}
        \item Lower layers: syntax
        \item Higher layers: semantics
        \item Learns linguistic structure automatically
    \end{itemize}

    \item \textbf{Revolutionary Impact}
    \begin{itemize}
        \item Established pre-training as standard
        \item Democratized NLP research
        \item Foundation for modern LLMs
    \end{itemize}
\end{enumerate}
\end{frame}

% ============================================================
% REFERENCES
% ============================================================
\begin{frame}{References 📚}
\textbf{Essential Papers:}

\vspace{0.5em}

\begin{itemize}
    \item \textbf{Devlin et al. (2019)} - "BERT: Pre-training of Deep Bidirectional Transformers for Language Understanding"
    \begin{itemize}
        \item The original BERT paper
        \item Introduced MLM and NSP
    \end{itemize}

    \vspace{0.3em}

    \item \textbf{Tenney et al. (2019)} - "BERT Rediscovers the Classical NLP Pipeline"
    \begin{itemize}
        \item Analysis of what BERT learns
        \item Layer-wise linguistic properties
    \end{itemize}

    \vspace{0.3em}

    \item \textbf{Clark et al. (2019)} - "What Does BERT Look At? An Analysis of BERT's Attention"
    \begin{itemize}
        \item Understanding BERT's attention patterns
    \end{itemize}
\end{itemize}

\vspace{0.5em}

\textbf{Tutorials:}
\begin{itemize}
    \item HuggingFace Course: Chapter 1 (Transformer Models)
    \item Jay Alammar: "The Illustrated BERT, ELMo, and co."
    \item BertViz: Interactive attention visualization
\end{itemize}
\end{frame}

% ============================================================
% QUESTIONS
% ============================================================
\begin{frame}{Questions? 🙋}
\begin{center}
\LARGE
\textbf{Discussion Time}
\end{center}

\vspace{2em}

\textbf{Topics for discussion:}
\begin{itemize}
    \item Masked Language Modeling
    \item Pre-training vs fine-tuning
    \item Contextual embeddings
    \item BERT architecture details
    \item Implementation questions
\end{itemize}

\vspace{2em}

\begin{center}
\Large
Thank you! 🙏

\vspace{1em}

Next: BERT Variants and Improvements!
\end{center}
\end{frame}

\end{document}
